\subsection{Digital Twin Definition}

\begin{frame}[more-space]{Summary: Digital Twin Definition}
  {\normalfont\small
  \renewcommand{\arraystretch}{1.35}
  \setlength{\tabcolsep}{4pt}
  \begin{tabular}{@{}p{0.23\textwidth}p{0.73\textwidth}@{}}
    \toprule
    \textbf{Element} & \textbf{Our thermal digital twin}\\
    \midrule
    \textbf{Model} & \mdseries
      Planned: PCE vs PI-PCE with monotonicity constraints; degrees 1 / 2\\
    \midrule
    \textbf{Parameters} & \mdseries
      Surrogate coordinates: CPU power $P$, fan signal $F$, inlet temperature $T$\\
    \midrule
    \textbf{Inputs} & \mdseries
      Surrogate: $(P,T,F)$; controller disturbance: inlet temperature\\
    \midrule
    \textbf{Outputs} & \mdseries
      Surrogate: maximum interval-mean CPU temperature; controller: proposed $(P,F)$\\
    \midrule
    \textbf{Measurement operator} & \mdseries
      18 CPU sensors, every 5 min; $\mathcal H=\max_{j=1,\ldots,18}\overline{T}_j$ over an approximately steady period\\
    \midrule
    \textbf{Observables} & \mdseries
      19 approximately steady temperature targets\\
    \midrule
    \textbf{Quantity of interest} & \mdseries
      CPU / fan choice within a validated domain and error bound; hardware protects peaks\\
    \bottomrule
  \end{tabular}
  }
\end{frame}
